%%%%%%%%%%%%%%%%%%%%%%%%%%%%%%%%%%%%%%%%%%%%%%%%%%%%%%%%%%%%%%%%%%%%%%%%%%%%%%%%%%%

\blankpage
\pagebreak
\thispagestyle{empty}
\addcontentsline{toc}{part}{\bfseries Portugal}
\mbox{}
\begin{center}

\vspace{3cm}

{\brabo\HUGE Portugal}
\end{center}
\pagebreak

%%%%%%%%%%%%%%%%%%%%%%%%%%%%%%%%%%%%%%%%%%%%%%%%%%%%%%%%%%%%%%%%%%%%%%%%%%%%%%%%%%%

\blankpage
\pagebreak
\thispagestyle{empty}
\addcontentsline{toc}{part}{Ao redor de Lisboa}
\mbox{}
\begin{center}

\vspace{3cm}

% {\formularlight\small PARTE 1}

% \medskip

{\brabo\HUGE\itshape Ao redor de Lisboa}
\end{center}
\pagebreak

%\part[Drama barroco e tragédia]{Drama barroco\\ e tragédia}

%%%%%%%%%%%%%%%%%%%%%%%%%%%%%%%%%%%%%%%%%%%%%%%%%%%%%%%%%%%%%%%%%%%%%%%%%%%%%%%%%%%

\chapter{Alfama}

Alfama é um dos bairros mais antigos de Lisboa e preserva um traçado de ruas estreitas, escadarias e pequenas praças que atravessaram diferentes períodos da história da cidade. A área concentra parte importante do patrimônio medieval e do fado.

Vale caminhar sem pressa entre a Sé de Lisboa, as ruas da antiga Judiaria, a Igreja de Santo António, o Museu do Fado e os miradouros de Santa Luzia e das Portas do Sol. O bairro é especialmente interessante para observar a sobreposição de arquiteturas e formas de ocupação urbana.

\chapter{Mouraria}

A Mouraria é um dos bairros históricos associados à formação do fado e à diversidade cultural de Lisboa. O bairro se desenvolveu junto às antigas muralhas da cidade e recebeu diferentes comunidades ao longo dos séculos.

\chapter{Baixa e Praça do Comércio}

A Baixa foi reconstruída após o terremoto de 1755 segundo um plano urbano regular, associado ao Marquês de Pombal. Suas ruas retas e quarteirões organizados contrastam com o traçado irregular de Alfama e Mouraria.

A Praça do Comércio, aberta para o Tejo, ocupa o local do antigo Terreiro do Paço e funciona como um dos principais espaços monumentais da cidade. A partir dali, vale caminhar pela Rua Augusta até o Rossio e o Chiado.

\chapter{Chiado e Bairro Alto}

O Chiado concentra livrarias, teatros, cafés, lojas históricas e instituições culturais. A região está associada à vida literária de Lisboa e à presença de escritores e artistas desde o século \textsc{xix}.

O Bairro Alto, por sua vez, preserva uma malha urbana antiga e uma intensa vida noturna. A região também possui casas de fado, bares, restaurantes e pequenos espaços culturais.

\chapter{Livrarias e cultura do livro}

Lisboa possui uma forte tradição livreira, com livrarias históricas e espaços contemporâneos dedicados à literatura.

Entre os endereços mais interessantes estão a Bertrand do Chiado, fundada em 1732, a Sá da Costa e a Ler Devagar. Também vale observar pequenas livrarias independentes e alfarrabistas, especialmente nas regiões do Chiado, Baixa e Bairro Alto.

\chapter{Gulbenkian}

A Fundação Calouste Gulbenkian é um dos principais centros culturais de Lisboa. O complexo reúne museu, jardins, biblioteca e programação de música, exposições e artes performativas.

A coleção do Museu Gulbenkian inclui arte europeia e importantes conjuntos de arte islâmica, oriental e de outras regiões. Para esta viagem, é especialmente interessante pela possibilidade de observar relações artísticas entre o Mediterrâneo, o Oriente e a Europa.

\chapter{Azulejos}

A azulejaria é uma das formas artísticas mais características da cultura portuguesa. Os azulejos aparecem em igrejas, palácios, estações, fachadas e espaços domésticos, assumindo funções decorativas e narrativas.

O Museu Nacional do Azulejo é a principal instituição dedicada à história dessa tradição, mas atualmente está temporariamente fechado para obras de restauro. Vale verificar sua situação antes da viagem.

\chapter{Distrito de Belém}

\begin{itemize}
\item Mosteiro dos Jerónimos: É uma das grandes obras do manuelino, construído no início do século \textsc{xvi}, e Patrimônio Mundial da \textsc{unesco} desde 1983.

\item Igreja de Santa Maria de Belém: Fica dentro do conjunto dos Jerónimos.

\item Padrão dos Descobrimentos: Vale mais pelo contexto histórico do que pela arquitetura.

\item Torre de Belém: Pode ser vista por fora, não precisa necessariamente entrar.

\item Museu Nacional de Arte Antiga: A coleção inclui arte portuguesa e europeia e também objetos de origem africana e asiática. Mas atenção: atualmente o Museu Nacional de Arte Antiga aparece como temporariamente fechado para obras.
\end{itemize}

\chapter{Évora}

\begin{itemize}
\item Templo Romano: É um dos grandes símbolos da cidade e data do século \textsc{i}, no período de Augusto. E uma observação interessante: ele é chamado popularmente de ``Templo de Diana'', mas hoje se sabe que essa identificação não é correta; estava ligado ao culto imperial.

\item Sé de Évora: Vale entrar, não apenas ver por fora. É uma das grandes construções medievais portuguesas e fica praticamente ao lado do templo romano --- uma sobreposição histórica muito bonita.

\item Igreja de São Francisco \& Capela dos Ossos: É famosa, mas não é só aquela atração macabra de Instagram. A igreja reúne elementos góticos, mouriscos, manuelinos, renascentistas, barrocos e azulejos; a Capela dos Ossos é do século \textsc{xvii}.

\item Praça do Giraldo: Eu não gastaria muito tempo na praça, mas começaria ou terminaria o passeio ali. É o centro natural para caminhar pela cidade histórica.

\item Universidade de Évora: Vale passar pela antiga Universidade, fundada no século \textsc{xvi}. É interessante justamente pensar Évora como centro de produção e circulação de conhecimento, e não apenas como cidade monumental. A promoção turística também a inclui entre os pontos que merecem uma visita.

\item Judiaria: O que existe é o local que provavelmente correspondeu à antiga sinagoga, na área da Travessa do Barão. Há documentação medieval, referências à ``Sinagoga Velha'' e à ``Sinagoga Grande'', e estudos que identificam ali uma possível sala de orações. Também existe uma marca interpretada como mezuzá em um portal.
\end{itemize}